\documentclass[12pt,a4paper]{article}

% 中文支持
\usepackage{ctex}

% 数学公式
\usepackage{amsmath}
\usepackage{amssymb}
\usepackage{amsthm}

% 图片支持
\usepackage{graphicx}

% 表格支持
\usepackage{booktabs}
\usepackage{array}

% 页面设置
\usepackage{geometry}
\geometry{left=2.5cm,right=2.5cm,top=2.5cm,bottom=2.5cm}

% 超链接
\usepackage{hyperref}
\hypersetup{
    colorlinks=true,
    linkcolor=blue,
    filecolor=magenta,
    urlcolor=cyan,
}

% 代码高亮（可选）
\usepackage{listings}
\lstset{
    basicstyle=\ttfamily,
    breaklines=true,
    frame=single
}

% 标题信息
\title{\textbf{中文学术论文示例模板}}
\author{作者姓名}
\date{\today}

\begin{document}

\maketitle

\begin{abstract}
这是一篇使用 \LaTeX 编写的中文学术论文示例。本文展示了如何使用 \LaTeX 排版中文文档，包括数学公式、图片、表格等常用元素。本模板适用于中文学术论文写作，提供了基本的排版框架和常用功能。
\end{abstract}

\section{引言}

\LaTeX 是一个高质量的文档排版系统，特别适合用于学术论文、技术文档等专业文档的编写。使用 \texttt{ctex} 宏包，我们可以方便地编写中文文档。

本文档演示了以下功能：
\begin{itemize}
    \item 中文排版
    \item 数学公式
    \item 图片插入
    \item 表格制作
    \item 代码展示
\end{itemize}

\section{数学公式}

\subsection{行内公式}

在文本中可以插入行内公式，例如 $E = mc^2$ 是爱因斯坦的质能方程，$a^2 + b^2 = c^2$ 是勾股定理。

\subsection{行间公式}

复杂的公式可以单独成行展示：

\begin{equation}
    \int_{-\infty}^{\infty} e^{-x^2} dx = \sqrt{\pi}
\end{equation}

多行公式可以使用 \texttt{align} 环境：

\begin{align}
    f(x) &= x^2 + 2x + 1 \\
         &= (x + 1)^2
\end{align}

矩阵表示：

\begin{equation}
    A = \begin{pmatrix}
        a_{11} & a_{12} & a_{13} \\
        a_{21} & a_{22} & a_{23} \\
        a_{31} & a_{32} & a_{33}
    \end{pmatrix}
\end{equation}

\section{图片}

图片是论文中的重要元素。下面展示如何插入图片：

\begin{figure}[htbp]
    \centering
    % 取消注释以下行以插入实际图片
    % \includegraphics[width=0.6\textwidth]{images/example.png}
    
    % 占位符：如果没有图片，可以使用空白框
    \fbox{\parbox{0.6\textwidth}{\centering \vspace{3cm} 在此插入图片 \\ \texttt{images/example.png} \vspace{3cm}}}
    
    \caption{图片示例}
    \label{fig:example}
\end{figure}

在文中可以引用图片，如图~\ref{fig:example}~所示。

\section{表格}

表格是展示数据的有效方式。使用 \texttt{booktabs} 宏包可以制作专业的三线表：

\begin{table}[htbp]
    \centering
    \caption{实验数据示例}
    \label{tab:example}
    \begin{tabular}{lccr}
        \toprule
        \textbf{项目} & \textbf{数值1} & \textbf{数值2} & \textbf{结果} \\
        \midrule
        实验A & 12.3 & 45.6 & 57.9 \\
        实验B & 23.4 & 56.7 & 80.1 \\
        实验C & 34.5 & 67.8 & 102.3 \\
        \bottomrule
    \end{tabular}
\end{table}

表~\ref{tab:example}~展示了实验数据的整理方式。

\section{代码展示}

如果需要在论文中展示代码，可以使用 \texttt{listings} 宏包：

\begin{lstlisting}[language=Python, caption=Python代码示例]
def hello_world():
    print("Hello, World!")
    return True

if __name__ == "__main__":
    hello_world()
\end{lstlisting}

\section{定理和证明}

\newtheorem{theorem}{定理}
\newtheorem{lemma}{引理}
\newtheorem{definition}{定义}

\begin{definition}
连续函数：设函数 $f(x)$ 在点 $x_0$ 的某个邻域内有定义，如果 $\lim_{x \to x_0} f(x) = f(x_0)$，则称 $f(x)$ 在点 $x_0$ 处连续。
\end{definition}

\begin{theorem}
设 $f(x)$ 和 $g(x)$ 在点 $x_0$ 处连续，则 $f(x) + g(x)$ 和 $f(x) \cdot g(x)$ 在点 $x_0$ 处也连续。
\end{theorem}

\begin{proof}
根据连续函数的定义和极限的运算法则可以直接证明。
\end{proof}

\section{结论}

本文展示了使用 \LaTeX 编写中文学术论文的基本方法，包括文字排版、数学公式、图片、表格等常用元素。使用本模板，研究者可以专注于内容创作，而不必过多担心排版问题。

\LaTeX 的优势在于：
\begin{enumerate}
    \item 专业的排版效果
    \item 强大的数学公式支持
    \item 自动化的交叉引用
    \item 便于版本控制
\end{enumerate}

\section*{致谢}

感谢开源社区对 \LaTeX 和相关工具的贡献。

% 参考文献（可选）
\begin{thebibliography}{99}
\bibitem{latex}
Lamport, Leslie. \textit{\LaTeX: A Document Preparation System}. Addison-Wesley, 1994.

\bibitem{ctex}
CTeX 开发组. \textit{CTeX 宏集手册}. 2021.
\end{thebibliography}

\end{document}
